\section{Safety 与 Security：统一视角下的目标函数}

讲者明确区分但又统一了两类目标。`AI safety` 关注系统是否对外部环境造成伤害；`AI security` 关注系统是否被外部攻击者利用、篡改、劫持。Agent 语境下两者高度耦合，因为攻击者可以利用安全缺陷触发安全事故，也可以利用安全机制漏洞突破系统约束。

\begin{knowledgebox}{课程中的定义}
\textbf{Safety}：防止系统对人类、组织和环境造成不希望的外部危害。\\
\textbf{Security}：防止系统被恶意主体破坏机密性、完整性、可用性（CIA）并被利用执行恶意行为。\\
在 Agent 场景里，二者共同构成 ``resilient alignment''：即使处于对抗环境，系统也能维持目标一致性。
\end{knowledgebox}

从工程角度，可以把风险粗略写成：
$$
R \approx P(\text{attack succeeds}) \times \text{Impact} \times \text{Exposure}
$$
其中 Exposure 由权限范围、可调用工具、外部环境耦合深度决定。Agent 的 Exposure 通常远大于纯文本模型，因此即便攻击成功率不变，整体风险也会显著上升。

\begin{warningbox}{常见误区}
把 ``模型回答更礼貌'' 当作系统更安全。现实里，很多高危攻击并不依赖有害文本，而依赖工具链、身份上下文、执行路径和权限配置。
\end{warningbox}

\subsection{本章小结}
Safety 与 Security 在 Agent 时代必须一起做。安全目标不是单指标优化，而是围绕行为后果、系统韧性和攻击成本的综合约束。

